\documentclass[letterpaper,11pt]{article}

\usepackage{latexsym}
\usepackage[empty]{fullpage}
\usepackage{titlesec}
\usepackage{marvosym}
\usepackage[usenames,dvipsnames]{color}
\usepackage{verbatim}
\usepackage{enumitem}
\usepackage[hidelinks]{hyperref}
\usepackage{fancyhdr}
\usepackage[english]{babel}
\usepackage{tabularx}
\input{glyphtounicode}

\pagestyle{fancy}
\fancyhf{}
\fancyfoot{}
\renewcommand{\headrulewidth}{0pt}
\renewcommand{\footrulewidth}{0pt}

\addtolength{\oddsidemargin}{-0.5in}
\addtolength{\evensidemargin}{-0.5in}
\addtolength{\textwidth}{1in}
\addtolength{\topmargin}{-.5in}
\addtolength{\textheight}{1.0in}

\urlstyle{same}

\raggedbottom
\raggedright
\setlength{\tabcolsep}{0in}

\titleformat{\section}{
  \vspace{-4pt}\scshape\raggedright\large
}{}{0em}{}[\color{black}\titlerule \vspace{-5pt}]

\pdfgentounicode=1

\newcommand{\resumeItem}[1]{
  \item\small{
    {#1 \vspace{-2pt}}
  }
}

\newcommand{\resumeSubheading}[4]{
  \vspace{-2pt}\item
    \begin{tabular*}{0.97\textwidth}[t]{l@{\extracolsep{\fill}}r}
      \textbf{#1} & \textit{\small #4} \\
      \textit{\small#3} & \textit{\small #2} \\
    \end{tabular*}\vspace{-7pt}
}

\newcommand{\resumeProjectHeading}[2]{
    \item
    \begin{tabular*}{0.97\textwidth}{l@{\extracolsep{\fill}}r}
      \small#1 & #2 \\
    \end{tabular*}\vspace{-7pt}
}

\renewcommand{\labelitemi}{$\circ$}
\renewcommand\labelitemii{$\circ$}

\newcommand{\resumeSubHeadingListStart}{\begin{itemize}[leftmargin=0.15in, label={}]}
\newcommand{\resumeSubHeadingListEnd}{\end{itemize}}
\newcommand{\resumeItemListStart}{\begin{itemize}[leftmargin=0.2in]}
\newcommand{\resumeItemListEnd}{\end{itemize}\vspace{-5pt}}

\begin{document}

\begin{center}
    \textbf{\Huge \scshape Jodh Singh Nahal} \\ \vspace{1pt}
    \small 778-538-2351 $|$
    \href{mailto:nahal.jodhsingh@gmail.com}{\underline{nahal.jodhsingh@gmail.com}} $|$
    \href{https://www.linkedin.com/in/jodhsinghnahal/}{\underline{linkedin.com/in/jodhsinghnahal}} $|$
    \href{https://github.com/jodhsinghnahal}{\underline{github.com/jodhsinghnahal}}
\end{center}

\section{Education}
\resumeSubHeadingListStart
    \resumeSubheading
      {University of British Columbia}{Expected May 2028}
      {Bachelor of Applied Science in Computer Engineering}{Vancouver, BC}
      \resumeItemListStart
        \resumeItem{Presidential Scholars Award (\$60,000), Dean's Honour List, available 4 or 8 months, CGPA: 84.7\%.}
      \resumeItemListEnd
\resumeSubHeadingListEnd

\section{Technical Skills}
\begin{itemize}[leftmargin=0.15in, label={}]
\small{\item{
     \textbf{Languages}{: Python, MATLAB \& Simulink, C/C++, SQL} \\
     \textbf{Math \& Modeling}{: NumPy, PyTorch, LTspice, Simulink/Simscape, Waveform Analysis} \\
     \textbf{Tools}{: Linux, Git, Docker}
}}
\end{itemize}

\section{Experience}
\resumeSubHeadingListStart

    \resumeSubheading
      {Undergraduate Research Assistant}{May 2026 -- Sept 2026}
      {UBC Electrical and Computer Engineering --- Supervisor: Dr. Lutz Lampe}{Vancouver, BC}
      \resumeItemListStart
        \resumeItem{Modeling nonlinear SSHI interface circuits and verifying simulated voltage waveforms against theoretical models.}
        \resumeItem{Optimizing power extraction of piezoelectric energy harvesting interface circuits (SEH, SSHI, SECE) through LTspice and Simulink/Simscape simulation.}
        \resumeItem{Researching battery-less platforms for the Internet of Underwater Things via simulation of piezoelectric interface circuits.}
        \resumeItem{Designing embedded power-management PCB hardware in KiCad integrating AC--DC rectifiers, transformers, and switching circuitry.}
      \resumeItemListEnd

    \resumeSubheading
      {Software Test / Embedded Development Co-op}{May 2024 -- Aug 2025}
      {Xantrex}{Burnaby, BC}
      \resumeItemListStart
        \resumeItem{Achieved $7\times$ device capacity, 90\% lower CPU usage, and $\sim$2\,s lower latency after the C++ migration, using code profiling to eliminate bottlenecks.}
        \resumeItem{Wrote Python scripts to autogenerate scalable C++ object-oriented code from SQLite and Microsoft Access/Excel data, eliminating manual boilerplate.}
        \resumeItem{Developed an object-oriented, multi-threaded, non-blocking C++ firmware architecture on Linux-based Raspberry Pi gateways to parse RV-C CAN frames from STM32 devices and publish structured messages over MQTT.}
        \resumeItem{Replaced a $\sim$60k-line Python implementation with the C++ architecture.}
        \resumeItem{Extended firmware support to NMEA 2000 and a proprietary UART battery communication protocol.}
        \resumeItem{Built Python Hardware-in-the-Loop (HIL) CAN and UART simulators to run 100+ verification tests on the Raspberry Pi gateway, validating edge cases and stress-testing real-time performance.}
        \resumeItem{Troubleshot CAN/BLE communication faults on embedded inverter and charger platforms using oscilloscopes, logic analyzers, PCAN-USB adapters, and BLE sniffers.}
        \resumeItem{Reduced build times by $5\times$ by implementing CI/CD cross-compilation pipelines with Docker, Jenkins, and CMake.}
      \resumeItemListEnd

    \resumeSubheading
      {Hardware Verification Engineering Intern}{Sept 2026 -- Present}
      {AMD}{Vancouver, BC}
      \resumeItemListStart
        \resumeItem{Writing SystemVerilog Assertions (SVA) to cross-check legal responses across IPs.}
      \resumeItemListEnd

\resumeSubHeadingListEnd

\section{Projects}
\resumeSubHeadingListStart

    \resumeProjectHeading
      {\textbf{Adversarial Robustness Evaluation} $|$ \emph{PyTorch, CNN, Computer Vision}}{}
      \resumeItemListStart
        \resumeItem{Built a modular attack framework with configurable hyperparameters to analyze perturbation level versus Attack Success Rate (ASR).}
        \resumeItem{Implemented FGSM and PGD adversarial attack pipelines in PyTorch to evaluate VGG-16 robustness.}
        \resumeItem{Trained a VGG-16 CNN on CIFAR-10 to establish baseline and per-class accuracy.}
        \resumeItem{Authored tools to visualize adversarial noise and its impact on decision boundaries.}
      \resumeItemListEnd

    \resumeProjectHeading
      {\textbf{Behaviour Cloning for Robotic Control} $|$ \emph{PyTorch, MuJoCo, NumPy, dm\_control}}{}
      \resumeItemListStart
        \resumeItem{Improved reaching precision with a custom distance-weighted loss function.}
        \resumeItem{Designed a 2-layer DNN in PyTorch mapping robotic state observations to joint torque predictions for a simulated Reacher task.}
      \resumeItemListEnd

    \resumeProjectHeading
      {\textbf{Image Compression Quadtree (Toqutree)} $|$ \emph{C++, STL, Image Processing}}{}
      \resumeItemListStart
        \resumeItem{Implemented 2D prefix sum arrays for O(1) average color and hue histogram lookups over arbitrary regions.}
      \resumeItemListEnd

\resumeSubHeadingListEnd

\end{document}
